  \subsubsection{同じ値を取る区間を群としてまとめる}\label{節：ガウス型の群数列}
    \bookterm{floor-function}{床関数}は,中身が同じ二つの連続整数の間にある間,同じ値を取る.
    例えば,正整数$k,j$に対して,
    \[
      \lfloor\sqrt k\rfloor=j
      \iff j^2\le k\le(j+1)^2-1
    \]
    だから,値$j$が$2j+1$項続く.
    このようなまとまりを一つの群として数えるのが,群\bookterm{sequence}{数列}の考え方である.
    実際に$\lfloor\sqrt k\rfloor$は,
    \[
      \underbrace{1,1,1}_{3\text{個}},\,
      \underbrace{2,2,2,2,2}_{5\text{個}},\,
      \underbrace{3,3,3,3,3,3,3}_{7\text{個}},\,\ldots
    \]
    と並ぶ.

    \bookterm{difference}{階差}にこの形が現れたら,\sectionref{節：階差型}と同じように足す.
    \textbf{群の境界と,最後の群のどこまでを足すか}を区別しよう.
    似た形でも,和の上端が$n$か$n-1$かで境界が変わる.

    \noindent \bookblocklink{example-gauss-role-2}{problem}{answer}{【例題】}\par\nobreak
    \begin{enumerate}[label=(\arabic*),parsep=3pt,beginpenalty=10000,format=\bookitemlink{example-gauss-role-2}{problem}{answer}]
      \item \bookterm{sequence}{数列}を次のように定める.\bookterm{general}{一般項}を求めよ.
            \[
              a_1=2,\qquad
              a_{n+1}=a_n+\lfloor\sqrt{n+2}\rfloor
              \quad(n\ge1).
            \]
            必要ならば,\,$m=\lfloor\sqrt{n+1}\rfloor$を用いてよい.
    \end{enumerate}

    \noindent \bookblocklink{example-gauss-role-2}{hint}{problem}{【方針】}\par\nobreak
    \begin{enumerate}[label=(\arabic*),parsep=3pt,beginpenalty=10000,format=\bookitemlink{example-gauss-role-2}{hint}{problem}]
      \item \bookterm{difference}{階差}を足し,\,$\lfloor\sqrt k\rfloor$の和にします.
            \,$k=1,2$で値がともに1であることを使い,和の範囲をそろえましょう.
    \end{enumerate}

    \noindent \bookblocklink{example-gauss-role-2}{answer}{problem}{【解答】}\par\nobreak
    \begin{enumerate}[label=(\arabic*),parsep=3pt,beginpenalty=10000,format=\bookitemlink{example-gauss-role-2}{answer}{problem}]
      \item \,$m=\lfloor\sqrt{n+1}\rfloor$とおくと,
            \[
              a_n=m(n+2)-\frac{m(m+1)(2m+1)}6.
            \]
    \end{enumerate}

    \noindent \bookblocklink{example-gauss-role-2}{solution}{problem}{【解法】}\par\nobreak
    \begin{enumerate}[label=(\arabic*),parsep=3pt,beginpenalty=10000,format=\bookitemlink{example-gauss-role-2}{solution}{problem}]
      \item \bookterm{difference}{階差}を足すと,
            \[
              a_n=2+\sum_{k=3}^{n+1}\lfloor\sqrt k\rfloor
                  =\sum_{k=1}^{n+1}\lfloor\sqrt k\rfloor.
            \]
            \,$n=1$では最初の和を空の和として0とする.
            \,$N=n+1,\ m=\lfloor\sqrt N\rfloor$とおく.
            第$m-1$群までは全項,第$m$群は$k=m^2$から$N$までの$N-m^2+1$項を足すので,
            \begin{align*}
              a_n&=\sum_{j=1}^{m-1}j(2j+1)+m(N-m^2+1)\\
                 &=\frac{(m-1)m(4m+1)}6\\
                 &\quad+m(N-m^2+1)\\
                 &=m(N+1)-\frac{m(m+1)(2m+1)}6.
            \end{align*}
            \,$N=n+1$を戻せば\bookterm{general}{一般項}を得る.\,$m=1$でも最初の和は0となり,初項と一致する.

            \noindent \textbf{【別解】}\\
            \,$\lfloor\sqrt k\rfloor$は,\,$j^2\le k$を満たす正整数$j$の個数でもある.
            同じ組$(j,k)$を$j$ごとに数えると,
            \begin{align*}
              a_n&=\sum_{j=1}^{m}(N-j^2+1)\\
                 &=m(N+1)-\frac{m(m+1)(2m+1)}6.
            \end{align*}
            群ごとにまとめる方法と,何を数えているかを読み直す方法がつながっている.
    \end{enumerate}
